\documentclass[12pt,a4paper]{article}
%per compilare il template usare questo sotto e commentare l'altro:
%\usepackage{../stile}
\usepackage{../../../template/stile}


%Titolo documento
\newcommand{\titoloDocumento}{Norme di Progetto}

%Prima data di creazione del documento
\newcommand{\dataCreazione}{30 novembre 2015}

%Inserite la versione attuale del documento
\newcommand{\versione}{1.0.0}

%Stato in cui si trova il documento: Formale solo all'atto di consegna
\newcommand{\stato}{Informale}

%Uso del documento
\newcommand{\uso}{Interno}

\rhead{\titoloDocumento}
\lfoot{Versione: \versione}
\title{\titoloDocumento}

\begin{document}
\begin{titlepage}
\begin{center}
\today \\
\vspace{1cm}
\begin{Huge}
\textbf{\nomeTeam} \\
\end{Huge}
\textbf{\CapitolatoL} \\
\vspace{1cm}
\includegraphics[scale=0.5]{\logoGrande}
\vspace{1cm}

\HRule \\[0.4cm]
\begin{Huge}
{\huge \bfseries \titoloDocumento}\\[0.4cm]
\end{Huge}
\HRule \\[1cm]
\vfill

\begin{table}[h]
\begin{center}
\begin{tabular}{r | l}
\multicolumn{2}{c}{\textbf{Informazioni sul documento}}\\
\midrule
\textbf{Nome Documento}	&	\titoloDocumento	\\
\textbf{Versione}	&	\versione	\\
\textbf{Stato}	&	\emph{\stato}	\\
\textbf{Uso}	&	\emph{\uso}	\\
\textbf{Data Creazione}	&	\dataCreazione	\\
\textbf{Data Ultima Modifica}	&	\today	\\
\textbf{Redazione}	&	\NDC	\\
\ &	\AVE	\\
\ &	\AVI	\\
\textbf{Verifica}	&	\IB	\\
\ & \TP \\
\textbf{Approvazione}	&	\AB	\\
\textbf{Lista Distribuzione}	&	\nomeTeam	\\

\end{tabular}
\end{center}
\end{table}

\end{center}
\end{titlepage}
\newpage

\Large{\textbf{Registro delle modifiche}}\\
\normalsize

\begin{table}[h]
\begin{center}

\begin{tabular}{p{0.12\textwidth} p{0.25\textwidth} p{0.18\textwidth} p{0.4\textwidth}}
\toprule
\textbf{Versione}	&	\textbf{Autore}	&	\textbf{Data}	&	\textbf{Descrizione}\\
\midrule
\midrule
1.0.0 & \NDC & 2015-11-30 & Template di base \\
\bottomrule
\end{tabular}
\caption{Versionamento del documento}
\label{tabVers1}
\end{center}
\end{table}
\newpage

\tableofcontents
\newpage

\listoftables
\listoffigures
\newpage

% NOTA non so dove mettere l'uso di git

\section{Introduzione}
Questo documento definisce le norme che i membri del gruppo \nomeTeam{} adotteranno nello svolgimento del progetto \CapitolatoL. Tutti i membri sono tenuti a leggere il documento attentamente e a seguirene le norme.

\subsection{Scopo}
Il documento si propone come proposito di garantire l’uniformità di tutto il materiale prodotto, migliorarene l’efficienza e ridurrene gli errori.

\subsection{Descrizione}
Verranno definite norme riguardanti:

\begin{itemize}
  \item interazioni tra membri del gruppo;
  \item stesura documenti, convenzioni stilistiche e tipografiche;
  \item modalità di lavoro durante tutte le fasi del progetto;
  \item ambiente di lavoro.
\end{itemize}

\subsection{Glossario}
Al fine di evitare ogni ambiguità di linguaggio e massimizzare la comprensione dei
documenti, i termini tecnici, di dominio, gli acronimi e le parole che necessitano chiarimenti sono riportate nel documento \emph{Glossario\_v1.2.0}. Ogni occorrenza di vocaboli presenti nel \emph{Glossario} è marcata da una G maiuscola in pedice (e.g. pedice\textsubscript{G}).

\newpage

\section{Processi di Sviluppo} %ISO 12207 5

\subsection{Scopo}

\subsection{Descrizione}

\subsection{Analisi dei Requisiti} %ISO 12207 5.3.2

\newpage

\section{Processi di Supporto} %ISO 12207 6

\subsection{Scopo}
Definisce norme per lo sviluppo e il mantenimento della documentazione prodotta durante il ciclo di vita del software. Inoltre definisce metodi per il controllo della qualità, di verifica e validazione di tali documenti.


\subsection{Descrizione}
Verranno definite norme riguardanti:

\begin{itemize}
  \item utilizzo e l'accesso ai documenti;
  \item templates da utilizzare, l'uniformità di linguaggio, le convenzioni stilistiche e tipografiche;
  \item metodi di verifica e approvazione dei documenti;
\end{itemize}

\subsection{Documentazione} %ISO 12207 6.1

\subsubsection{Template}
Per agevolare la redazione della documentazione è stato creato un template \LaTeX{} contenente tutte le impostazioni stilistiche e grafiche citate in questo documento. Tale modello si può trovare nella repository\textsubscript{G} in \texttt{doc/template}.

\subsubsection{Struttura dei documenti}

\subsubsubsection{Prima pagina}
Ogni documento deve avere una prima pagina che contiene le seguenti informazioni sul documento:

\begin{itemize}
  \item nome del gruppo;
  \item nome del progetto;
  \item logo del gruppo;
  \item titolo del documento;
  \item versione del documento;
  \item stato del documento (informale\textsubscript{G} o formale\textsubscript{G})
  \item uso del documento (interno\textsubscript{G} o esterno\textsubscript{G});
  \item data di creazione del documento;
  \item data di dultima modifica del documento;
  \item nome e cognome dei redattori del documento;
  \item nome e cognome dei verificatori del documento;
  \item nome e cognome del responsabile che approva il documento;
  \item lista di distribuzione del documento;     
\end{itemize}

\subsubsubsection{Diario delle modifiche}
Il diario delle modifiche del documento segue la prima pagina. La tabella deve essere ordinata per data in ordine decrescente, in modo che la prima riga corrisponda alla versione attuale del documento. Ogni riga del diario delle modifiche deve contenere:

\begin{itemize}
  \item versione del documento dopo la modifica;
  \item nome e cognome dell'autore della modifica;
  \item data della modifica;
  \item titolo del documento;
  \item una breve descrizione delle modifiche svolte.
\end{itemize}

\subsubsubsection{Indice}
Al diario delle modifiche deve sempre seguire un indice delle sezioni e sottosezioni.

\subsubsubsection{Elenco tabelle e figure}
All'indice delle sezioni e sottosezioni può seguire un elenco delle tabelle e delle figure. Nel caso non siano presenti figure o tabelle i rispettivi indici devono essere
omessi.

\subsubsubsection{Introduzione}
\paragraph{Scopo}
\paragraph{Descrizione}

\subsubsubsection{Sezioni}
\paragraph{Scopo}
\paragraph{Descrizione}

\subsubsubsection{Formattazione delle pagine}
L’intestazione di ogni pagina deve contenere:

\begin{itemize}
  \item logo del gruppo;
  \item nome del documento.
\end{itemize}

Il piè di pagina deve contenere:
\begin{itemize}
  \item versione del documento;
  \item nome dell'università e anno accademico corrente;
  \item pagina corrente nel formato N di T dove N è il numero di pagina corrente e T è il numero di pagine totali;
  \item email del gruppo;
  \item licenza di distribuzione del documento.
\end{itemize}

\subsection{Linguaggio}
Nei documenti alcuni termini devono essere interpretati in modo analogo ai termini inglesi descritti in RFC 2119\footnote{http://www.ietf.org/rfc/rfc2119.txt}:

\begin{itemize}
  \item i termini \emph{deve}, \emph{è richiesto} e sinonimi stretti sono da intendersi con lo stesso significato di \emph{MUST};
  \item i termini \emph{non deve}, \emph{è richiesto che non} e sinonimi stretti sono da intendersi con lo stesso significato di \emph{MUST NOT};
  \item i  termini \emph{dovrebbe}, \emph{si raccomanda}, \emph{è preferibile} e sinonimi stretti sono da intendersi con lo stesso significato di \emph{SHOULD}
  \item i termini \emph{non dovrebbe}, \emph{si raccomanda di non}, \emph{è preferibile che non} e sinonimi stretti sono da intendersi con lo stesso significato di \emph{SHOULD NOT};
  \item i termini \emph{può}, \emph{opzionalmente} e sinonimi stretti sono da intendersi con lo stesso significato di \emph{MAY}.
\end{itemize}

\subsection{Configurazione} %ISO 12207 6.2

\subsection{Verifica} %ISO 12207 6.4

\subsection{Validazione} %ISO 12207 6.5

\subsection{Risoluzione dei Problemi} %ISO 12207 6.8

\newpage

\section{Processi di Organizzazione} %ISO 12207 7.1

\subsection{Scopo} %ISO 12207 7.1.1

\subsection{Descrizione} %ISO 12207 7.1.1

\subsection{Pianificazione} %ISO 12207 7.1.2
%task ruoli e responsabilità

\subsection{Esecuzione e controllo} %ISO 12207 7.1.3
%come viene monitorata l'esecuzione della pianificazione: DA DEFINIRE (per esempio chiedendo ogni tanto come sta andando il lavoro)

\subsection{Infrastruttura} %ISO 12207 7.2
%asana per i task, github per le issues e bug e risoluzione problemi, projectlibre per la pianificazione con gantt e costi

\newpage


\end{document}
